% ==========================================================================
% Preliminaries + Method for: Plücker Rotary Addressing for Test-Time-Training
% Camera Conditioning (Idea 2: qk_rope_cam)
% Written to be \input{} into an ICLR-style template.
% ==========================================================================

\newcommand{\ours}{PRA\xspace}   % Plücker Rotary Addressing
\newcommand{\R}{\mathbb{R}}
\newcommand{\ip}[2]{\left\langle #1, #2 \right\rangle}
\newcommand{\silu}{\mathrm{silu}}
\newcommand{\NS}{\mathrm{NS}}

\section{Preliminaries}
\label{sec:prelim}

We first fix notation for the multi-view synthesis setting
(\S\ref{sec:prelim-nvs}), review how rotary embeddings achieve
\emph{relative} conditioning in attention (\S\ref{sec:prelim-rope}), and then
describe the large-chunk test-time-training (LaCT) layer that our method
targets (\S\ref{sec:prelim-lact}). Finally we recall Plücker coordinates,
the ray parameterization on which our encoding is built
(\S\ref{sec:prelim-plucker}).

\subsection{Multi-view novel view synthesis with LVSM}
\label{sec:prelim-nvs}

Given $V_{\mathrm{in}}$ posed input images
$\{(I_v, K_v, E_v)\}_{v=1}^{V_{\mathrm{in}}}$ with intrinsics
$K_v \in \R^{3\times3}$ and camera-to-world extrinsics
$E_v \in \mathrm{SE}(3)$, the task is to render the scene from
$V_{\mathrm{tg}}$ novel target poses. LVSM-style models
cast this as sequence modeling: every image is split into
$P$ non-overlapping patches, and each patch becomes a token. Input-view
tokens carry the patch RGB concatenated with a per-pixel ray map; target-view
tokens carry only the ray map. A transformer-like backbone processes the
concatenated sequence of $L = (V_{\mathrm{in}} + V_{\mathrm{tg}})\,P$ tokens,
and target tokens are linearly decoded to RGB patches.

Throughout, each token $i$ is associated with the viewing ray of its patch
center: the ray origin $\mathbf{o}_i \in \R^3$ (the camera center of the view
that produced the token) and the unit direction $\mathbf{d}_i \in \R^3$,
both expressed in a \emph{canonical scene frame} obtained by normalizing the
camera set (mean pose to identity, positions scaled into the unit box).
We use row-vector convention: features $x \in \R^{1 \times d}$ multiply
weights from the right, matching common implementations.

\subsection{Relative conditioning via rotary embeddings in attention}
\label{sec:prelim-rope}

Rotary position embeddings (RoPE) inject positions
into attention logits so that only \emph{relative} positions survive.
For a scalar position $t_i$, RoPE partitions the head dimension into 2D
pairs and rotates each pair by a position-proportional angle. Writing
$R(\theta) \in \mathrm{SO}(2)$ for the rotation of one pair and
$\theta_i = \omega t_i$,
\begin{equation}
\label{eq:rope-identity}
\ip{R(\omega t_j)\, q}{R(\omega t_i)\, k}
  \;=\;
\ip{q}{R\big(\omega (t_i - t_j)\big)\, k},
\end{equation}
i.e., rotations meet inside the inner product and compose into a function of
the \emph{difference} $t_i - t_j$ alone. Absolute inputs, relative
interactions: this identity is the entire mechanism.

Camera-pose extensions replace scalar positions by per-view transforms. GTA
and PRoPE apply block-diagonal matrix representations
of the camera transform to queries, keys, and values, so that attention
logits contain $q^\top G_j^\top G_i k$: for orthogonal (or appropriately
inverted projective) $G$, the product $G_j^\top G_i$ encodes the relative
pose between the views of tokens $i$ and $j$. These constructions are
attention-specific: they rely on the fact that every query--key interaction
in attention is \emph{exactly} a bilinear form. Whether an analogous
guarantee can be obtained for nonlinear fast-weight layers is the question
addressed in this work.

\subsection{Test-time training layers and LaCT}
\label{sec:prelim-lact}

A test-time-training (TTT) layer replaces the
attention operator with a small \emph{fast-weight} network $f_W$ whose
parameters $W$ are updated \emph{during the forward pass} by gradient
descent on a self-supervised regression loss, then queried. LaCT
scales this recipe with (i) a large SwiGLU
fast-weight network and (ii) \emph{large-chunk} updates: all key--value
pairs of the input views form a single chunk and contribute to one
(or few) gradient step(s), which keeps GPU utilization high.

Concretely, per attention head the LaCT layer used in LVSM maintains
$W = (W_0, W_2, W_1)$ with $W_0, W_2 \in \R^{d \times d_h}$,
$W_1 \in \R^{d_h \times d}$, and computes
\begin{equation}
\label{eq:swiglu-fw}
f_W(x) \;=\; \big(\, \silu(x W_0) \odot (x W_2) \,\big)\, W_1 ,
\end{equation}
where $\odot$ is the elementwise product. Queries, keys and values are
produced from the token stream by a learned projection
$(q, k, v) = \silu(x\, W_{qkv})$, and $q, k$ are $\ell_2$-normalized per
token. The layer executes two phases:

\textbf{Update (write).} For the chunk of input-view tokens
$\{(k_i, v_i)\}_{i \in \mathcal{I}}$, the fast weights take one gradient
step on the negative-inner-product regression loss
$\mathcal{L}(W) = -\sum_{i \in \mathcal{I}} \ell r_i \ip{v_i}{f_W(k_i)}$,
where $\ell r_i > 0$ are per-token learning rates predicted from the token
content by a small head. Writing $h(x) = \silu(x W_0) \odot (x W_2)$ for
the hidden activation, the resulting updates are outer products
\begin{align}
\Delta W_1 &\;\propto\; \textstyle\sum_{i} \big(\ell r_{1,i}\, h(k_i)\big)^{\!\top} v_i,
\label{eq:w1-update}\\
\Delta W_0 &\;\propto\; \textstyle\sum_{i} \big(\ell r_{0,i}\, k_i\big)^{\!\top} g_i,
\qquad
\Delta W_2 \;\propto\; \textstyle\sum_{i} \big(\ell r_{2,i}\, k_i\big)^{\!\top} u_i,
\label{eq:w0w2-update}
\end{align}
with $g_i, u_i \in \R^{1 \times d_h}$ the backpropagated signals of the
SwiGLU gate and hidden branches (functions of $k_i$, $v_i$, and the
pre-update weights). LaCT additionally orthogonalizes each update with a
few Newton--Schulz iterations (Muon) and renormalizes
weight columns to their initial norms; both operations preserve the
outer-product \emph{structure} of \eqref{eq:w1-update}--\eqref{eq:w0w2-update}
in the sense that the update remains a sum over tokens of rank-one terms in
the $k$- (resp.\ $h(k)$-) direction, reweighted across tokens.

\textbf{Apply (read).} Every token $j$ (input and target) queries the
updated weights: $o_j = f_{W + \Delta W}(q_j)$.

In the LaCT-LVSM architecture, camera information enters the network
\emph{only once}, through the ray map concatenated to the RGB input at
patchification. The TTT layer itself---the only module that exchanges
information across views---is pose-blind: neither the write
\eqref{eq:w1-update}--\eqref{eq:w0w2-update} nor the read has any explicit
notion of which view a token came from, or where its ray lives in space.

\subsection{Plücker coordinates}
\label{sec:prelim-plucker}

A ray $\{\mathbf{o} + t\,\mathbf{d} : t \ge 0\}$ with unit direction
$\mathbf{d}$ determines the \emph{Plücker coordinates}
\begin{equation}
\label{eq:plucker}
\pi \;=\; (\mathbf{d},\, \mathbf{m}) \in \R^6,
\qquad
\mathbf{m} \;=\; \mathbf{o} \times \mathbf{d}.
\end{equation}
The moment $\mathbf{m}$ is invariant to sliding $\mathbf{o}$ along the ray,
so $\pi$ identifies the underlying \emph{line} irrespective of
parameterization. Two geometric facts make Plücker coordinates a natural
token address for multi-view synthesis. First, $\pi$ separates rays: two
tokens see the world along the same line iff their coordinates agree.
Second, rays that intersect at a common scene point $\mathbf{p}$ satisfy
$\mathbf{m}_i = \mathbf{p} \times \mathbf{d}_i$, hence for two such rays
\begin{equation}
\label{eq:moment-bound}
\mathbf{m}_i - \mathbf{m}_j \;=\; \mathbf{p} \times (\mathbf{d}_i - \mathbf{d}_j),
\qquad
\|\mathbf{m}_i - \mathbf{m}_j\| \;\le\; \|\mathbf{p}\|\, \|\mathbf{d}_i - \mathbf{d}_j\|,
\end{equation}
i.e., rays observing the same 3D point from nearby directions are provably
close in Plücker space, with the moment discrepancy controlled by the scene
radius. This is the property that will let a memory addressed by $\pi$
associate content across views.

% ==========================================================================
\section{Method: Plücker Rotary Addressing for Fast Weights}
\label{sec:method}

Our goal is to make every LaCT layer camera-aware while (i) preserving the
relative-encoding guarantees that make RoPE-style methods effective in
attention, (ii) exploiting---rather than fighting---the structure of the
fast-weight update, and (iii) leaving the architecture, the compiled update
kernel, and the optimization dynamics untouched. We first show that the
LaCT layer, despite its nonlinearity, reads its newly written memory
exclusively through inner products (\S\ref{sec:method-memory}); this single
observation transplants the RoPE identity \eqref{eq:rope-identity} to fast
weights. We then construct the encoding (\S\ref{sec:method-pra}), analyze
what is and is not relative (\S\ref{sec:method-analysis}), and discuss
stability and implementation (\S\ref{sec:method-impl}).

\subsection{Fast weights are inner-product-addressed memories}
\label{sec:method-memory}

Consider one LaCT layer operating in the single-chunk regime of
\S\ref{sec:prelim-lact}: all input tokens are written in one gradient step
taken at the pre-update weights $W^0 = (W_0^0, W_2^0, W_1^0)$, after which
every token reads. Expanding the read of token $j$ to first order in the
updates \eqref{eq:w1-update}--\eqref{eq:w0w2-update} gives
\begin{equation}
\label{eq:readout}
o_j
\;=\;
\underbrace{h^0(q_j)\, W_1^0}_{\text{init readout}}
\;+\;
\underbrace{\sum_{i \in \mathcal{I}} \ell r_{1,i}\,
   \ip{h^0(q_j)}{h^0(k_i)}\, v_i}_{\text{value retrieval}}
\;+\;
\underbrace{\sum_{i \in \mathcal{I}}
   \ip{q_j}{k_i}\; c_{ij}}_{\text{gate/hidden corrections}}
\;+\; O(\Delta W^2),
\end{equation}
where $c_{ij} \in \R^{1 \times d}$ collects the (token-dependent)
coefficients arising from $\Delta W_0$ and $\Delta W_2$ passing through the
SwiGLU. The derivation is elementary: every update is a sum of outer
products $a_i^\top b_i$ with $a_i \in \{k_i,\, h^0(k_i)\}$
(Eqs.~\ref{eq:w1-update}--\ref{eq:w0w2-update}), and a row vector meeting an
outer product contracts through an inner product,
$x\, (a^\top b) = \ip{x}{a}\, b$.

Equation~\eqref{eq:readout} is the structural fact our method rests on, so
we state its consequence explicitly:

\begin{lemma}[Inner-product addressing]
\label{lem:addressing}
In the single-chunk regime, every interaction between a query $q_j$ and the
content written by a key $k_i$---through any of $\Delta W_0, \Delta W_1,
\Delta W_2$, at any order of the expansion---enters exclusively through the
inner products $\ip{q_j}{k_i}$ or $\ip{h(q_j)}{h(k_i)}$. The pre-update
weights $W^0$ are the only path by which $q_j$ affects the output not
through an inner product with some $k_i$.
\end{lemma}

The lemma says that a LaCT layer is an associative memory with a
\emph{kernel} readout: writes deposit rank-one payloads indexed by keys, and
reads retrieve them weighted by query--key inner products---exactly the
algebraic situation \eqref{eq:rope-identity} was designed for. Attention
needed no such analysis because its logits are bilinear by definition; here
the same property holds for the \emph{update-induced} part of a nonlinear
layer, which is what test-time training adds to the network.

\subsection{Plücker rotary addressing}
\label{sec:method-pra}

\ours{} equips every token with camera-geometry-dependent rotations applied
to $q$ and $k$ in every TTT layer. Let $\pi_i = (\mathbf{d}_i, \mathbf{m}_i)
\in \R^6$ be the Plücker coordinates \eqref{eq:plucker} of token $i$'s
patch-center ray in the canonical scene frame, and let
$\omega_1 < \dots < \omega_F$ be a geometric ladder of frequencies. Define
$6F$ phases per token,
\begin{equation}
\label{eq:phases}
\theta_i^{(a,f)} \;=\; \gamma_{a,f}\;\omega_f\; \pi_i^{(a)},
\qquad a = 1, \dots, 6,\quad f = 1, \dots, F,
\end{equation}
where $\pi_i^{(a)}$ is the $a$-th Plücker coordinate and $\gamma_{a,f}$ are
learnable per-coordinate, per-frequency gains initialized to $1$. Collect
the phases into the block-diagonal rotation
\begin{equation}
\label{eq:block-rot}
G_i \;=\;
\mathrm{blockdiag}\Big(
R\big(\theta_i^{(1,1)}\big), \dots, R\big(\theta_i^{(6,F)}\big),\,
I_{d - 12F}
\Big) \;\in\; \mathrm{SO}(d),
\end{equation}
which rotates the first $12F$ feature dimensions in independent 2D planes
and leaves the remaining $d - 12F$ dimensions untouched as a pure content
channel. After the layer's usual $\ell_2$ normalization of queries and
keys, \ours{} applies
\begin{equation}
\label{eq:pra-apply}
\tilde q_j \;=\; q_j\, G_{j}^{\top},
\qquad
\tilde k_i \;=\; k_i\, G_{i}^{\top},
\end{equation}
and the layer proceeds \emph{unchanged}: the same update kernel
\eqref{eq:w1-update}--\eqref{eq:w0w2-update} writes $\tilde k_i$ against
$v_i$, and reads with $\tilde q_j$. Values, learning rates, Muon, and
weight normalization are untouched. In our instantiation $d = 256$,
$F = 16$, so $12F = 192$ dimensions are rotated and $64$ remain
content-only; the frequencies span
$\omega_f \in \pi \cdot [0.5,\, 16]$, geometrically spaced, matched to the
canonical scene scale in which camera positions lie in the unit box (so the
lowest frequency disambiguates coarse ray orientation while the highest
resolves sub-degree, sub-percent-of-scene ray offsets).

The computation is a single fused elementwise operation--- $O(LF)$
multiply-adds on top of an $O(L d\, d_h)$ layer---and adds $6F = 96$
learnable scalars per layer. No new modules, losses, or kernel changes are
introduced.

\subsection{What is relative, and what is not}
\label{sec:method-analysis}

\paragraph{The update-induced kernel is exactly relative.}
Substituting \eqref{eq:pra-apply} into the retrieval terms of
Lemma~\ref{lem:addressing} and using orthogonality of $G$,
\begin{equation}
\label{eq:relative-kernel}
\ip{\tilde q_j}{\tilde k_i}
\;=\;
q_j\, G_j^\top G_i\, k_i^\top
\;=\;
\ip{q_j}{\,k_i\,\big(G_j^\top G_i\big)^{\!\top}}
,
\qquad
G_j^\top G_i
=
\mathrm{blockdiag}\Big( R\big(\theta_i^{(a,f)} - \theta_j^{(a,f)}\big) \Big)_{a,f},
\end{equation}
so every phase enters the memory readout only through the
\emph{differences} $\theta_i^{(a,f)} - \theta_j^{(a,f)} =
\gamma_{a,f}\,\omega_f\,(\pi_i^{(a)} - \pi_j^{(a)})$: the write/read kernel
depends on camera geometry exclusively through the relative Plücker offset
$\pi_i - \pi_j$ between the writing and reading rays. The same argument
applies verbatim to the $\ip{h(\tilde q_j)}{h(\tilde k_i)}$ channel
whenever the hidden nonlinearity is evaluated on rotated inputs whose
interaction is mediated by the $\Delta W$ outer products
(Lemma~\ref{lem:addressing}). In particular, by
\eqref{eq:moment-bound}, two tokens in different views observing the same
scene point have phase offsets bounded by scene radius times their
directional parallax across \emph{all} frequencies simultaneously---the
memory retrieves content written by rays through the same 3D region, which
is precisely the inductive bias multi-view synthesis needs.

\paragraph{The initial-weight residue.}
By Lemma~\ref{lem:addressing} the only non-relative path is the pre-update
weights: terms of the form $\tilde q_j W_0^0$ (and analogously through
$W_2^0$, $W_1^0$) expose the \emph{absolute} phase $G_j$. Two properties
render this residue benign. First, at initialization the columns of
$W_0^0, W_2^0$ are isotropic Gaussian, whose distribution is invariant under
the orthogonal maps $G_j$; the encoding therefore leaves the forward
statistics of the network exactly unchanged at the start of training, and
the slow weights co-adapt to the rotated input distribution thereafter.
Second, the fast weight is a \emph{per-scene} state: because all rays are
expressed in a canonical frame fixed by the input camera set, ``absolute''
phases are already relative to the scene gauge---across scenes, no global
frame is shared for the slow weights to overfit to. This argument, which
has no analogue in attention (where the shared slow weights mediate
\emph{all} token interactions), is particular to test-time training and is
the reason an orthogonal, RoPE-style construction suffices where attention
required projective machinery.

\paragraph{Multi-view consistency of the memory.}
It is instructive to read \eqref{eq:relative-kernel} as a statement about
what the fast weight stores. With \ours{}, the rank-one payloads
$\tilde k_i^\top(\cdot)$ deposited by the update live on phase-modulated
key directions; a query retrieves a superposition in which each
contribution is weighted by a product of cosines of relative Plücker
offsets across the frequency ladder,
\begin{equation}
\label{eq:kernel-shape}
\kappa(\pi_i, \pi_j) \;=\; \textstyle\sum_{c} \prod_{(a,f) \in c}
\cos\!\big(\gamma_{a,f}\,\omega_f\,(\pi_i^{(a)} - \pi_j^{(a)})\big)\,(\dots),
\end{equation}
a translation-invariant kernel on Plücker space, sharply peaked at
$\pi_i = \pi_j$ and decaying with ray disparity at a rate set by the learned
gains. The fast weight thus behaves as a memory over the scene's
\emph{ray space}: input views write appearance along their rays; a target
pixel reads along its own ray and aggregates content from geometrically
consistent writes. Cross-view association---previously the burden of
content similarity alone---is provided by geometry at every layer.

\subsection{Stability and implementation}
\label{sec:method-impl}

\paragraph{Orthogonality is load-bearing.}
$G_i \in \mathrm{SO}(d)$ preserves the $\ell_2$ normalization of $q$ and
$k$ exactly, so the calibration of the inner-loop loss, the per-token
learning rates, the Muon orthogonalization, and the weight-norm constraint
are all numerically identical to the baseline. This is what allows \ours{}
to be inserted into the compiled update kernel's inputs without touching
the kernel: the encoding commutes with everything downstream of the
projections. (Non-orthogonal alternatives---e.g., projective PRoPE
transforms---rescale tokens and act as uncontrolled perturbations of the
write strengths; we evaluate such a variant in our experiments.)

\paragraph{Where to inject.}
\ours{} is applied \emph{after} the $\ell_2$ normalization, in every TTT
layer, to $q$ and $k$ only. Applying it before normalization would let the
normalization partially undo the rotation only if norms varied per block,
which they do not; the post-norm placement makes the invariance exact and
keeps the implementation a two-line change. Values are deliberately left
untouched: the value pathway admits its own exact relative treatment,
which is orthogonal to (and composable with) \ours{}.

\paragraph{Learnable gains.}
The gains $\gamma_{a,f}$ (initialized to 1) let training reweight---or
switch off---individual coordinate/frequency channels. This subsumes
frequency-selection heuristics: if a moment coordinate at a high frequency
aliases for a given scene-scale distribution, its gain can decay toward
zero, recovering the baseline behavior for that channel. At
$\gamma \equiv 0$ the layer is exactly the pose-blind baseline, so the
baseline is contained in \ours{}'s hypothesis class.

\paragraph{Complexity and integration.}
Per layer, \ours{} costs one elementwise rotation of $2 \times 12F$
dimensions per token ($<0.1\%$ of layer FLOPs) and $6F$ parameters; the
phases are computed once per forward pass from the patch-center rays
already produced by the input pipeline, and are shared across layers up to
the per-layer gains. Training requires no schedule changes, no auxiliary
losses, and no modification of the LaCT kernel, checkpoint format, or
inference-time reconstruct/render split.
